\begin{table}[h!]
\centering
\caption{Kategorien und Zielgrößen aus der CTQ}
\label{table:zielgroeßen}
\begin{tabular}{|l|l|}
\hline
Kategorie  & Zielgrößen  \\
\hline
\textbf{Geometrie}  & \textbf{Durchmesser, Höhe, Ausbreitungsfaktor}\\
Masse  & Gewicht, Gewichtsverlust beim Backen \\
Struktur  & Härte, Knusprigkeit, Bruchkraft \\
Optik  & Bräunungsgrad, Rissbildung, Gleichmäßigkeit der Kontur \\
Sensorik  & Süße, Buttergeschmack, Mundgefühl \\
\hline
\end{tabular}
\end{table}